\chapter{Il nucleo di CPU}
\label{cap:kernel-cpu}

Il nucleo di CPU non è la parte più veloce del progetto, ma è quella di cui
tutto il resto ha bisogno per essere giudicato: è il denominatore di ogni
speedup, ed è l'unica implementazione che gira ovunque, anche su una macchina
senza \code{nvcc}. Merita quindi lo stesso lavoro dei backend CUDA, e per le
stesse ragioni.

Il vincolo che se ne è dato il progetto è che sia C11 portabile, senza
intrinsics e senza direttive di vettorizzazione. Non è una rinuncia: è una
scelta che sposta il problema. Senza intrinsics l'unica leva sulla
vettorizzazione è la \emph{forma} del codice, che deve essere abbastanza
semplice e abbastanza esplicita da lasciare al compilatore una decisione facile.
Buona parte di questo capitolo riguarda proprio come si scrive un ciclo perché
un compilatore lo vettorizzi da sé.

\section{Tre backend, una sola idea}
\label{sec:tre-backend}

Tutti e tre i backend di CPU implementano lo schema~A di \cref{sec:schemi}:
cambia come il lavoro viene organizzato attorno a quell'ordine dei cicli, non
l'ordine. Si scelgono a tempo di compilazione con la variabile \code{KERNEL}
del \code{Makefile}, espongono la stessa interfaccia (\cref{app:interfaccia}) e
i binari coesistono, così un confronto è una sola campagna sugli stessi dati.

\begin{itemize}[nosep,left=0pt]
  \item \kernelname{scheme\_a} --- l'implementazione di riferimento, quella che
        corrisponde alla consegna;
  \item \kernelname{scheme\_a\_jblock} --- la stessa, con $X$ mattonellata in
        cache;
  \item \kernelname{scheme\_a\_jblock\_rb} --- come sopra, più il
        \emph{register blocking} sulle righe di $A$.
\end{itemize}

I tre non sono tre tentativi alternativi ma una \textbf{catena}: ciascuno
rimuove il collo di bottiglia che il precedente lascia scoperto. Il capitolo li
presenta in quest'ordine, e ogni sezione si apre con il limite che motiva il
backend successivo.

\section{Il backend di riferimento}

\subsection{Perché $k$ deve essere noto al compilatore}
\label{sec:specializzazione}

Nello schema~A i $k$ accumulatori sono variabili locali del ciclo su $j$, e
tutto il vantaggio dello schema dipende dal fatto che vivano nei
\textbf{registri}: è lì che il valore $A[i][j]$, letto una volta sola, viene
riusato $k$ volte.

Perché il compilatore possa tenerceli deve però sapere \emph{quanti} sono
mentre compila. Il motivo è concreto: i registri non sono indirizzabili a
runtime. Se il numero di accumulatori è noto solo quando il programma parte,
\code{acc} diventa per forza un array in memoria e ogni aggiornamento
\code{acc[c] += \ldots} un accesso a quella memoria --- esattamente il traffico
che lo schema~A serviva a evitare.

Il codice risolve generando \textbf{cinque funzioni specializzate}, una per
ciascuno dei $k$ richiesti dalla traccia, a partire da una sola definizione via
macro:

\begin{lstlisting}[language=C,caption={Generazione delle specializzazioni
(estratto di \file{src/kernel/scheme\_a.c}).}]
#define COLS_8(M)  COLS_6(M) M(6) M(7)
#define DECLARE_ACC(c) scalar_t acc##c = (scalar_t)0;
#define UPDATE_ACC(c)  acc##c += a * xrow[c];
#define STORE_ACC(c)   yrow[c] = acc##c;

/* per ogni riga i: */
COLS(DECLARE_ACC)
for (j = 0; j < n_loc; j++) {
    const scalar_t a = arow[j];          /* letto UNA volta */
    const scalar_t *restrict xrow = X_loc + (size_t)j * ldx;
    COLS(UPDATE_ACC)                     /* k aggiornamenti espliciti */
}
COLS(STORE_ACC)
\end{lstlisting}

Vale la pena leggere che cosa producono queste tre righe. \code{COLS\_8} è un
elenco di indici; applicarlo a \code{DECLARE\_ACC} genera otto dichiarazioni di
variabili scalari \emph{distinte} (\code{acc0}, \code{acc1}, \ldots), non un
array; applicarlo a \code{UPDATE\_ACC} genera otto aggiornamenti scritti per
esteso, senza alcun ciclo interno con estremo runtime. Il percorso caldo che il
compilatore vede non contiene quindi nessuna indicizzazione su $c$: contiene
$k$ istruzioni indipendenti su $k$ variabili indipendenti, che è la forma in
cui l'allocazione in registro e la vettorizzazione sono decisioni facili.

\subsection{Il fallback generico, e perché la traccia lo richiede}

La traccia chiede però che il codice funzioni per $k$ \emph{generico}, e cinque
specializzazioni non possono coprire tutti i valori possibili. Per ogni altro
$k$ interviene un \textbf{fallback generico}, che non tenta di tenere $k$
accumulatori --- sarebbero troppi, e finirebbero in memoria --- ma lavora a
gruppi di al più $\code{KB} = 32$ colonne per volta. Ogni gruppo richiede una
passata completa su $A$:
\begin{equation}
  \text{passaggi su } A \;=\; \Bigl\lceil \frac{k}{\code{KB}} \Bigr\rceil .
  \label{eq:passaggi}
\end{equation}

Per $k \le \code{KB}$ il ciclo esterno sui gruppi è \emph{degenere}: esiste nel
codice ma gira una volta sola, $A$ viene letta una volta sola e il fallback si
differenzia dalle specializzazioni soltanto per l'array di accumulatori invece
delle variabili distinte. Oltre quella soglia i passaggi diventano due, e con
essi raddoppia il traffico sia su $A$ sia su $X$.

La conseguenza va tenuta presente leggendo i risultati: le righe dei $k$ oltre
la soglia e quelle dei $k$ entro la soglia \textbf{non sono direttamente
confrontabili}, perché non stanno facendo la stessa quantità di lavoro di
memoria. È l'andamento a dente di sega già anticipato in \cref{sec:schemi}.

\subsection{Il microbenchmark che isola la specializzazione}

Che le specializzazioni servano davvero è un'ipotesi, e va misurata. Compilando
con \code{FORCE\_GENERIC\_K=1} le cinque funzioni non vengono nemmeno generate
e ogni $k$ passa dal fallback: le due build differiscono per una sola variabile,
coesistono come binari distinti e \code{kernel\_name()} le distingue nel CSV
(\kernelname{scheme\_a} contro \kernelname{scheme\_a\_generic}). Il confronto è
quindi un A/B sullo stesso problema e nella stessa sessione di misura, non due
esecuzioni accostate a posteriori.

La campagna include un \textbf{controllo di consistenza} che vale la pena
dichiarare, perché è ciò che rende il risultato credibile: sui $k$ oltre la
soglia entrambe le build passano dal fallback, quindi lì le due righe devono
coincidere entro il rumore. Se non coincidessero, la differenza misurata sugli
altri $k$ non sarebbe specializzazione ma rumore, e la campagna andrebbe
rifatta.

\section{Perché $\code{KB} = 32$}

Il valore copre tutti i $k$ del collaudo, ma questa da sola non è una ragione
sufficiente: sarebbe una costante scelta per far tornare i casi che interessano.
La ragione vera è un budget di registri.

Con $\code{KB}$ accumulatori il compilatore deve tenere $\code{KB}$ valori vivi
attraverso l'intero ciclo su $j$. Vettorizzati, occupano
$\code{KB}/W$ registri vettoriali, dove $W$ è il numero di scalari che entrano
in un registro --- quattro in doppia precisione su AVX2, otto su AVX-512 --- a
fronte dei 16 o 32 registri architetturali disponibili. Al crescere di
$\code{KB}$ il riuso di $A$ migliora, perché i passaggi della
\cref{eq:passaggi} diminuiscono, ma cresce la pressione sui registri fino allo
\emph{spill}, cioè fino al punto in cui il compilatore è costretto a
parcheggiare in memoria gli accumulatori che non stanno più nei registri --- e
a quel punto si perde più di quanto si guadagni.

Il valore $32$ è il compromesso che tiene i cinque $k$ richiesti in un
passaggio solo restando sotto la soglia di spill sulle architetture di
interesse. Che lo sia davvero anche sui $k$ grandi è una previsione, non un
dato, ed è uno degli oggetti della campagna C3 (\cref{sec:c3}).

\section{Ciò che il compilatore vede}

La build usa \code{-O3} con \code{-march=native} (oppure \code{-mcpu=native} su
aarch64: il \code{Makefile} prova quale delle due il compilatore accetta), più
\code{-Wall -Wextra -Wpedantic}. Nessun \code{\#pragma}, nessun intrinsic.

Perché questo basti, il codice deve rispondere da sé alle tre domande che un
compilatore si pone prima di vettorizzare un ciclo.

\begin{itemize}[nosep,left=0pt]
  \item \textbf{«Questi puntatori possono riferirsi alla stessa memoria?»}
        No: $A$, $X$ e $Y$ sono dichiarati \code{restrict}. La garanzia va però
        \emph{ripristinata} dentro \code{local\_gemm}, con copie locali dei
        puntatori: il contesto opaco tiene il puntatore ad $A$ come membro di
        una struttura, e senza quella copia i kernel specializzati perderebbero
        l'informazione di non aliasing che avevano quando $A$ era un parametro.
        Se l'aliasing fosse possibile, ogni scrittura su $Y$ potrebbe
        invalidare $A$ e il ciclo andrebbe eseguito in ordine.
  \item \textbf{«Le iterazioni dipendono l'una dall'altra?»} No: i $k$
        aggiornamenti \code{acc[c] += a * X[j][c]} riguardano accumulatori
        distinti e sono indipendenti fra loro. Essendo già srotolati dalle
        macro, si presentano come catene di moltiplicazione-accumulo
        indipendenti, che è la forma che una unità FMA può eseguire in
        parallelo.
  \item \textbf{«Gli accessi sono regolari?»} Sì, e a passo unitario: come
        osservato in \cref{sec:schemi}, sia la riga di $A$ sia la riga di $X$
        sono contigue in memoria.
\end{itemize}

\begin{daverificare}
Verificare sul server l'effettiva vettorizzazione del percorso caldo, per
esempio con \code{gcc -fopt-info-vec} oppure ispezionando l'assembly di
\code{kernel\_k8} e \code{kernel\_k32}. Il punto da documentare è se e a quale
$k$ compaiano spill degli accumulatori.
\end{daverificare}

\section{Il primo limite: la rilettura di $X$}
\label{sec:rilettura-x}

Lo schema~A legge $A$ una volta sola, ed è il suo pregio. Ma per ogni riga $i$
di $A$ percorre \textbf{tutta} $X$, e questo è il suo prezzo: il traffico in
lettura su $X$ vale
\begin{equation}
  \text{elementi di } X \text{ letti} \;=\; \mloc \cdot \nloc \cdot k,
  \label{eq:traffico-x}
\end{equation}
cioè $k$ volte il traffico su $A$.

Finché $X$ resta in una cache vicina al calcolo questo traffico non raggiunge
la memoria centrale e semplicemente non si vede. Quando $X$ non ci sta più,
viene ristreamata da un livello più lontano una volta per ogni riga di $A$, e
il tempo per elemento smette di essere governato dal calcolo per essere
governato dalla banda di quel livello.

La soglia è esplicita e dipende da entrambi i parametri del problema: il blocco
locale $X_{\text{loc}}$ occupa $\nloc \cdot k \cdot s$ byte. A parità di
matrice, quindi, \textbf{è $k$ che decide da che parte della soglia ci si
trova}, ed è per questo che l'effetto si manifesta come un calo dei GFLOPS al
crescere di $k$ --- cioè nella direzione opposta a quella che l'intensità
aritmetica della \cref{eq:intensita} lascerebbe prevedere. È la chiave di
lettura dei risultati di \cref{sec:c3}, ed è il limite che il backend
successivo rimuove.

\section{Secondo backend: $X$ mattonellata in cache}
\label{sec:jblock}

L'idea di \kernelname{scheme\_a\_jblock} è di non percorrere più tutta $X$ per
ogni riga di $A$, ma di prenderne un \textbf{tile} --- un blocco di righe
consecutive --- e di applicarlo a tutte le righe di $A$ prima di passare al
successivo:

\begin{lstlisting}[language=C,caption={Struttura di \file{src/kernel/scheme\_a\_jblock.c}.}]
for j0 (tile di X di BJ righe):
  for i (TUTTE le righe di A):
    acc = (primo tile) ? 0 : Y[i][*]
    for j nel tile:  a = A[i][j];  acc[c] += a * X[j][c]
    Y[i][*] = acc
\end{lstlisting}

Il tile viene caricato dalla memoria una volta e riusato $\mloc$ volte
restando in cache. $A$ continua a essere letta \emph{una volta sola}: ogni sua
riga viene semplicemente consumata a segmenti, uno per tile, invece che tutta
di seguito.

\paragraph{Il prezzo.} Y non può più vivere nei registri per tutta la durata
di una riga, perché ci si torna sopra a ogni tile: va riletta e riscritta ogni
volta. Il traffico aggiuntivo vale $2k$ scalari per riga e per tile, cioè un
rapporto $2k/\text{BJ}$ rispetto al traffico su $A$. Il tile va quindi scelto
\emph{grande} abbastanza da rendere questo termine trascurabile e
\emph{piccolo} abbastanza da stare in cache: le due richieste tirano in
direzioni opposte, e il punto di equilibrio non è deducibile a priori.

Per questo la dimensione del tile è un parametro di build,
\code{JBLOCK\_BYTES}, e non una costante scelta una volta per tutte: il codice
ne ricava il numero di righe per tile dividendo per $k \cdot s$, il default
è un valore che sta nella cache di secondo livello di qualunque core x86
recente, e quale sia il valore migliore è oggetto di misura.

\paragraph{Un principio già visto.} È esattamente la stessa mossa del backend
CUDA con tile in shared memory (\cref{cap:kernel-cuda}): portare un blocco di
$X$ in una memoria vicina al calcolo e riusarlo per molte righe prima di
passare al blocco successivo. Cambia il nome della memoria vicina --- lì la
shared memory, qui la cache --- non il ragionamento.

\section{Terzo backend: register blocking sulle righe}
\label{sec:jblock-rb}

Tolto il limite di banda, ne resta scoperto un altro, che si manifesta ai $k$
\emph{piccoli}: in tutti i backend visti finora ogni \code{acc[c]} è una catena
\emph{seriale} di moltiplicazioni-accumulo lungo $j$ --- ogni aggiornamento
deve attendere il precedente --- e l'unico parallelismo a livello di istruzione
disponibile è quello fra le $k$ colonne. Poiché una unità FMA ha una latenza di
diversi cicli e la CPU ne ha più di una, per tenerle occupate servono più
catene indipendenti di quante $k$ piccolo ne offra: il processore aspetta, e i
GFLOPS crescono con $k$ invece di essere costanti.

\kernelname{scheme\_a\_jblock\_rb} risolve lavorando su $\text{RB}$ righe di
$A$ per volta:

\begin{lstlisting}[language=C,caption={Struttura di \file{src/kernel/scheme\_a\_jblock\_rb.c}.}]
for j0 (tile di X, come in scheme_a_jblock):
  for i a passi di RB:
    acc[r][c] = (primo tile) ? 0 : Y[i+r][c]
    for j nel tile:
      x[c] = X[j][c]                          /* caricato UNA volta */
      for r: a = A[i+r][j];  acc[r][c] += a * x[c]   /* riusato RB volte */
    Y[i+r][c] = acc[r][c]
\end{lstlisting}

Le righe di $A$ sono indipendenti per costruzione --- nessun elemento di $Y$
dipende da un'altra riga --- quindi il numero di catene indipendenti passa da
$\lceil k/W \rceil$ a $\text{RB} \cdot \lceil k/W \rceil$. In più, ogni valore
di $X$ caricato viene riusato $\text{RB}$ volte direttamente dai registri,
quindi anche il traffico su $X$ --- già in cache grazie al tile --- si riduce
dello stesso fattore.

\paragraph{Chi sceglie RB.} Anche qui il limite sono i registri:
$\text{RB} \cdot \lceil k/W \rceil$ accumulatori più $\text{RB}$ valori di $A$
devono starci tutti, altrimenti si torna allo spill e si perde più di quanto si
guadagni. Il codice sceglie perciò $\text{RB}$ \emph{a tempo di compilazione e
per ogni $k$}, dal numero di registri vettoriali e dalla loro ampiezza
dell'architettura di destinazione (\code{VEC\_NREGS} e \code{VEC\_W}, fissati
dalle macro del compilatore), prendendo il più grande fra i valori ammessi che
soddisfi il vincolo; la variabile di build \code{RB\_ROWS} permette di forzarlo
per misurarne l'effetto. Le righe che avanzano quando $\mloc$ non è multiplo di
$\text{RB}$ sono chiuse dal kernel a una riga, che è il caso $\text{RB} = 1$
dello stesso schema.

\medskip

I tre backend rispondono quindi a tre limiti diversi e successivi: il traffico
su $A$, risolto dallo schema~A stesso; la banda verso $X$, risolta dal tile in
cache; il parallelismo di istruzione, risolto dal blocco di righe. Quanto
ciascuno dei tre passi valga davvero sulla macchina di misura è ciò che la
campagna riportata in \cref{cap:risultati} quantifica.